% ------------------------------------------------------------------------------------
\section{Introduction}
Phase~1 of this project built classical machine-learning forecasters of highway speed on the METR-LA network~\cite{li2018dcrnn} (207 loop detectors, five-minute readings, March--June 2012) and measured one specific gap: a per-sensor XGBoost model given hand-built \emph{one-hop} upstream/downstream neighbour features improved by 4.1\% at 15 minutes but only 1.8\% at 60 minutes. The Lighthill--Whitham--Richards (LWR) theory~\cite{lighthill1955kinematic,richards1956shock} explains why: congestion waves travel upstream at 15--20~km/h, so what reaches a sensor an hour from now is still several sensors away, beyond the reach of one-hop features.

Phase~2 asks whether a deep model that \emph{learns} to pass information several hops along the road graph closes that gap, as the project brief proposes: ``train a Graph Attention Network or diffusion graph model over the road graph, using METR-LA sensor series and a fixed adjacency matrix''. Everything that defined the Phase~1 comparison is kept fixed --- the outage masking, the chronological 70/10/20 split with one-hour gaps, the one-hour input window, the 15/30/60-minute horizons, the exact test samples and the evaluation code --- so every difference in the numbers comes from the model.

Our contributions are: (1)~a spatio-temporal graph attention network (ST-GAT) with separate attention over upstream and downstream neighbours, implemented in plain PyTorch; (2)~a controlled comparison against a no-graph network of identical size and a fixed-weight diffusion graph network, which isolates the value of the graph and of attention; (3)~a \emph{hybrid} model in which the graph network learns the residual of the Phase~1 XGBoost forecast, trained on leakage-free out-of-fold predictions; (4)~a physics-informed variant with a soft LWR conservation penalty; (5)~seven controlled ablations, three seeds and Diebold--Mariano significance tests~\cite{diebold1995dm}; and (6)~an analysis of where the learned attention points in congested versus free-flowing traffic.

% ------------------------------------------------------------------------------------
\section{Deep Learning Literature Review}
\paragraph{Sequence models} Recurrent networks were the first deep models applied to traffic series~\cite{lv2015traffic}; LSTM~\cite{hochreiter1997lstm} and GRU~\cite{cho2014gru} cells keep a learned memory of the recent past, and DCRNN reports that an FC-LSTM without a graph reaches 3.44/3.77/4.37~mph MAE on METR-LA~\cite{li2018dcrnn}, \emph{worse} than a well-featured tree model, because each sensor is treated alone.

\paragraph{Graph convolutions} Treating sensors as graph nodes was the decisive step. DCRNN~\cite{li2018dcrnn} replaces the matrix products of a GRU by \emph{bidirectional diffusion convolutions} over the road graph's forward and backward random-walk matrices; STGCN~\cite{yu2018stgcn} uses spectral graph convolutions~\cite{kipf2017gcn} with gated temporal convolutions; T-GCN~\cite{zhao2019tgcn} couples a GCN with a GRU; Graph WaveNet~\cite{wu2019graphwavenet} learns an adaptive adjacency. All of these weight neighbours by a \emph{fixed} (or slowly learned) matrix.

\paragraph{Attention} Graph Attention Networks~\cite{velickovic2018gat} compute the weight of each neighbour from the current node states, so the same road link can matter a lot in congestion and little in free flow; GATv2~\cite{brody2022gatv2} makes the attention function more expressive. ASTGCN~\cite{guo2019astgcn}, ST-GAT~\cite{zhang2019stgat} and GMAN~\cite{shi2021masked} apply spatial (and temporal) attention to traffic and report gains over fixed-weight graph models, largest at long horizons~\cite{jiang2022survey}.

\paragraph{Hybrid and physics-informed models} Stacked generalisation~\cite{wolpert1992stacked} combines learners by training a second model on the first one's out-of-sample predictions; we use it to let a graph network correct a strong per-sensor tree model. Physics-informed neural networks add the residual of a governing equation to the loss~\cite{raissi2019pinn}; in traffic, LWR and second-order models have been used this way for state estimation~\cite{huang2020pidl,shi2020physics}. We adopt a light version: a soft LWR conservation penalty on the forecasts.

\paragraph{Gap} Published graph models are compared mostly with weak baselines (historical average, FC-LSTM). Few studies measure the gain of a graph network over a \emph{strong, neighbour-aware} classical model on identical splits and samples, separate the value of the graph from the value of deep sequence modelling, or test significance. Phase~2 does all three.

% ------------------------------------------------------------------------------------
\section{Data, Preprocessing and Regularisation}
\subsection{From the Phase 1 table to graph sequences}
The input of one example is a tensor $\mathbf{X}_t\in\mathbb{R}^{12\times 207\times 7}$ --- the last hour (steps $t-11,\dots,t$) of every sensor, with seven features per sensor and step --- and the target is $\mathbf{Y}_t\in\mathbb{R}^{207\times3}$, the speeds 15, 30 and 60 minutes later. The features are: (1)~speed, z-scored per sensor with \emph{training-period} mean and standard deviation; (2)~a mask (1 = observed, 0 = outage); (3--6)~sine/cosine of time of day and day of week; (7)~the sensor's historical-average speed for that time-of-week slot, computed from training weeks only. The 8.1\% of readings that are exactly 0~mph (offline sensors, Phase~1) are set to 0 \emph{after} scaling --- i.e.\ to the sensor's mean --- and flagged by the mask channel, so the network is told the value is missing instead of being told traffic stopped. Unknown targets are excluded from the loss and from every metric (masked MAE). Training, validation and test samples are exactly the Phase~1 sample indices (23{,}967 / 3{,}403 / 6{,}831 five-minute slots).

\subsection{The graph}
We use the fixed DCRNN adjacency $A_{ij}=\exp(-d_{ij}^2/\sigma^2)$ (thresholded at 0.1) built from road distances, which we re-derived and verified in Phase~1: 1{,}515 directed links, $A_{ij}>0$ meaning $j$ is downstream of $i$. Direction is kept: each sensor receives messages from its \emph{downstream} neighbours ($A_{ij}>0$) and, through a separate set of attention heads, from its \emph{upstream} neighbours ($A_{ji}>0$), each with a self-loop.

\subsection{Regularisation}
The models have 40--60 thousand parameters, against 4.6~million training targets per horizon, but consecutive windows are almost identical, so the effective sample size is far smaller. We regularise with: dropout~\cite{srivastava2014dropout} (0.1) on node states and on attention coefficients; decoupled weight decay (AdamW~\cite{loshchilov2019adamw}, $10^{-4}$); layer normalisation~\cite{ba2016layernorm} and residual connections in every graph layer; gradient-norm clipping at 5; a learning rate halved when validation error stalls for two epochs; and early stopping on validation masked MAE (patience 5). Each epoch draws 8{,}000 random training windows (a third of the training slots); this lowers the cost per epoch on CPU and adds sampling noise without changing the training distribution.
